% =====================================================================
% SECTION II — DROP-IN REPLACEMENT
% Incorporates resolutions of review issues #1 (exactness/ADF), #2
% (ridge, not pitchfork), #3 (coordinates and invariance), #6 (rho^2),
% #8 (validity typing).
%
% PREAMBLE ADDITIONS REQUIRED (after \usepackage{amsthm}):
%   \newtheorem{proposition}{Proposition}
%   \newtheorem{lemma}{Lemma}
%   \newcommand{\vd}{v_{d}}          % degeneracy direction
%   \newcommand{\rhosq}{\rho^{2}}    % discriminating SNR^2
%
% NUMBERS: every \manifest{...} placeholder is populated from
% results/manifest.tex generated by the analysis pipeline (T2--T11).
% Define \newcommand{\manifest}[1]{\textbf{[#1]}} until then.
%
% BIBLIOGRAPHY ADDITIONS: see comment block at end of file.
% =====================================================================

\section{Formalism: Exact Trajectories and Approximate Filters}
\label{sec:formalism}

\subsection{The S3 agent pair}

We define an S3 agent as a pair ($\rho$, $\Pi$), where $\rho$ is a
time-evolving probability measure over a parameter space $\Theta$ and
$\Pi$ is the structure that updates it: the prior together with the
likelihood of the data accumulated so far. For gravitational-wave
inference on a compact binary coalescence,
$\Theta = \{\Mchirp, q, \chieff, \Lambdatilde, \Dl\}$, and the natural
ordering variable is the frequency cutoff $f$, monotonic in time as the
signal chirps across the detection band. We write $\rho_f$ for the
posterior given data below cutoff $f$.

\subsection{Specialisation to GW inference}

For stationary Gaussian noise with one-sided power spectral density
$S_n(f)$, the log-likelihood at cutoff $f$ is
\begin{equation}
\ln \mathcal{L}(\theta \mid d, f) = -\frac12 \int_{\fmin}^{f}
\frac{|d(f') - h(f'; \theta)|^2}{S_n(f')}\,df' + \text{const},
\label{eq:likelihood}
\end{equation}
and the Fisher information matrix~{\cite{cutler1994fisher,vallisneri2008fisher}}
is
\begin{equation}
\Gamma_{ij}(f) = \int_{\fmin}^{f}
\frac{\partial_i h(f';\theta)\,\partial_j h(f';\theta)^{*} + \text{c.c.}}
{2\,S_n(f')}\,df',
\label{eq:fisher}
\end{equation}
evaluated at a stated expansion point (\secref{validity}). The
increment
\begin{equation}
\frac{d\Gamma_{ij}}{df} =
\frac{\partial_i h\,\partial_j h^{*} + \text{c.c.}}{2\,S_n(f)}
\label{eq:fisher_flow}
\end{equation}
is a deterministic flow on the cone of positive-semidefinite matrices,
driven entirely by the waveform model and the noise weighting.

\subsection{Exactness, and the two objects of study}
\label{sec:exactness}

Before introducing dynamics, we fix what can and cannot evolve. For
stationary Gaussian noise the likelihood factorises over disjoint
frequency bands, so Bayesian updating over those bands is a product of
commuting factors:

\begin{proposition}[Order independence of exact inference]
\label{prop:exact}
Let $\{B_k\}_{k=1}^{N}$ be a partition of $[\fmin,\fmax]$ into disjoint
bands with likelihood factors $\mathcal{L}_k(\theta)$. Then for any
permutation $\sigma$,
$\pi(\theta)\prod_k \mathcal{L}_{\sigma(k)}(\theta)
= \pi(\theta)\prod_k \mathcal{L}_{k}(\theta)$:
the exact posterior at $\fmax$ is independent of the order in which the
bands are assimilated.
\end{proposition}

Proposition~\ref{prop:exact} forbids any claim that the \emph{exact}
posterior depends on the path of information accumulation. Two
well-defined objects survive it, and they are the two subjects of this
paper.

\paragraph{Object 1: the exact posterior trajectory.}
The one-parameter family $\{\rho_f : \fmin \le f \le \fmax\}$, where
each $\rho_f$ is the exact posterior given data below $f$, is
well defined, and the frequency ordering of this family is physically
privileged: it is the time ordering experienced by any real-time
observer of a chirping binary, since data above $f$ does not yet exist
when the signal instantaneous frequency is $f$. The geometry of this
family --- how the local information metric reorients as $f$ grows
(\secref{commutator}) --- is an exact, observer-independent property of
the inference problem. Every endpoint statement is anchored by
Proposition~\ref{prop:exact}: $\rho_{\fmax}$ equals the batch
posterior.

\paragraph{Object 2: the single-pass Gaussian filter and its bias.}
Practical low-latency analyses cannot afford batch sampling and instead
propagate a Gaussian summary. We formalise the scheme as
assumed-density filtering (ADF)~{\cite{minka2001ep,opper1998online}} with
a Laplace projection: maintain
$q_k = \mathcal{N}(m_k, P_k)$; on receipt of band $B_{k+1}$, linearise
$h$ about $m_k$, perform the conjugate (Gauss--Newton) update, and
re-Gaussianise. Each projection discards non-Gaussian structure that
would otherwise interact with subsequent bands, so --- unlike the exact
update --- \emph{the single-pass filter is order dependent}, and its
terminal estimate carries a deterministic bias relative to the batch
posterior. Characterising that bias, and showing that it is governed by
the commutator structure of \eqnref{commutator_def}, is the second
subject of this paper (\secref{dynamic}, Appendix~{\ref{sec:appendix}}).
We emphasise the division of labour: nothing in Object~1 is an
approximation, and nothing in Object~2 is a property of exact
inference.

\subsection{Coordinates and invariance}
\label{sec:coordinates}

Angles, eigenvectors, and Frobenius norms in parameter space are
meaningful only relative to a metric, and the raw parameters carry
heterogeneous units: a Euclidean angle computed with $\Mchirp$ in solar
masses changes if $\Mchirp$ is expressed in seconds. We therefore fix,
once, the dimensionless analysis coordinates
\begin{equation}
x = \bigl(\ln \Mchirp,\; q,\; \chieff,\; \Lambdatilde/100,\; \ln \Dl\bigr),
\label{eq:coords}
\end{equation}
in which all eigendecompositions, projectors, and matrix norms in this
paper are computed; Fisher matrices and displacement vectors transform
by the Jacobian of \eqnref{coords} in the standard way. Robustness of
every headline quantity to two alternative conventions (a
symmetric-mass-ratio variant and a prior-range whitening) is reported
alongside the primary values (\secref{dynamic}).

Three classes of quantity require no convention at all, because the
Fisher matrix is the pullback of the waveform-space inner product
$\inner{a}{b} = 4\,\mathrm{Re}\!\int a b^{*}/S_n\,df$ onto parameter
space: (i) likelihood scalars such as
$\Delta\ln\mathcal{L} = \tfrac12\,\Delta x^{\top}\Gamma\,\Delta x$;
(ii) the exact mismatch
$\inner{\delta h}{\delta h}$ between two waveforms; and (iii) the
cumulative discriminating signal-to-noise between two parameter points,
\begin{equation}
\rhosq(f) \;\equiv\; \Delta x^{\top}\,\Gamma(f)\,\Delta x
\;\simeq\; \inner{\,h(\theta_1)-h(\theta_2)\,}{\,h(\theta_1)-h(\theta_2)\,}_{<f},
\label{eq:rhosq}
\end{equation}
where the right-hand side is the exact (convention-free) quantity and
the quadratic form its second-order approximation; we report both and
use the exact form wherever the displacement is not small
(\secref{validity}). Wherever possible, predictions in this paper are
stated in terms of these invariants.

\subsection{The projector, the commutator, and the misalignment $C(f)$}
\label{sec:commutator}

Let $\lambda_1(f) \ge \cdots \ge \lambda_5(f)$ and $v_1(f), \ldots,
v_5(f)$ be the eigenvalues and unit eigenvectors of $\Gamma(f)$ in the
coordinates of \eqnref{coords}. Define the projector onto the
best-constrained direction,
\begin{equation}
\Pi(f) = v_1(f)\,v_1(f)^{\top},
\label{eq:projector}
\end{equation}
and the commutator
\begin{equation}
\comm{\Gamma}{\Pi} \equiv \Gamma\Pi - \Pi\Gamma
= \lambda_1 v_1 v_1^{\top} - v_1 v_1^{\top}\Gamma .
\label{eq:commutator_def}
\end{equation}
The commutator vanishes iff $\Gamma$ has no off-diagonal structure
connecting $v_1$ to its orthogonal complement at that cutoff; it is the
infinitesimal generator of the eigenframe rotation of the cumulative
Fisher matrix. We quantify the rotation by the fractional misalignment
\begin{equation}
C(f) = \frac{\norm{\Gamma(f) - \lambda_1(f)\,\Pi(f)}_F}
{\norm{\Gamma(f)}_F},
\label{eq:cross_coupling}
\end{equation}
computed in the declared coordinates. We further denote by
$\vd(f) \equiv v_5(f)$ the \emph{degeneracy direction}: the eigenvector
of the smallest eigenvalue. For the binary-neutron-star systems
considered here, $\vd$ lies almost entirely in the
$(q, \chieff, \Lambdatilde)$ block at all cutoffs
[\manifest{T6: max component of $\vd$ outside the block}], reflecting
the well-known mass-ratio--spin--tidal degeneracy~{\cite{chatziioannou2020bns}}.
(Earlier drafts and some of the literature label this direction $v_2$
from analyses restricted to a low-dimensional subspace; we use $\vd$
throughout to avoid ambiguity about the ambient dimension.)

\subsection{The degeneracy ridge and prior lensing}
\label{sec:ridge}

At low cutoffs the likelihood \eqnref{likelihood} is nearly constant
along $\vd$: extended families of $(q, \chieff, \Lambdatilde)$
combinations produce waveforms whose accumulated discriminating
signal-to-noise \eqnref{rhosq} is below unity. The posterior in this
subspace is therefore a \emph{degeneracy ridge} --- a continuous,
elongated, prior-bounded structure --- rather than a collection of
isolated modes. We make no claim in this paper that the ridge fragments
into distinct maxima at any cutoff; whether it does is an open question
for time-resolved analysis (\secref{open}), and we deliberately avoid
bifurcation vocabulary until that analysis exists.

The ridge has an immediate, testable consequence: it determines how the
posterior responds to changes of prior.

\begin{lemma}[Prior lensing]
\label{lem:lensing}
Let the posterior be $p \propto \mathcal{L}\,\pi$ with likelihood
precision $\Gamma$ at the mean, and perturb the log-prior by a smooth
$\delta(\theta)$. To leading order in the Laplace expansion the
posterior mean shifts by
\begin{equation}
\Delta x_{*} \approx \Gamma^{-1}\,\nabla\delta
= \sum_{i=1}^{5} \lambda_i^{-1}\,(v_i^{\top}\nabla\delta)\,v_i ,
\label{eq:lensing}
\end{equation}
which for any perturbation with generic overlap
($v_d^{\top}\nabla\delta \neq 0$) is dominated by the degeneracy
direction: $\Delta x_{*} \parallel \vd$ up to corrections of order
$\lambda_d/\lambda_{i\neq d}$.
\end{lemma}

The inverse Fisher matrix acts as a lens that focuses arbitrary prior
perturbations into the degenerate subspace. For a hard truncation of
the prior (the case realised by the two published GW170817 spin priors)
the smooth-perturbation magnitude in \eqnref{eq:lensing} is not
controlled --- the mean shift saturates at the prior boundary --- but
the \emph{direction} of the shift remains that of the ridge, since
truncation removes probability mass from one end of an elongated
structure. Lemma~\ref{lem:lensing} therefore yields a directional
prediction, tested in \secref{static}: the displacement between the
posterior means of two prior-conditioned analyses of the same data
aligns with $\vd$, with alignment assessed against the null
distribution of a random direction in the relevant subspace.

\subsection{Domain of validity of the Fisher description}
\label{sec:validity}

The matrix $\Gamma(f)$ appears in this paper in three logically
distinct roles with different validity domains. Conflating them is the
standard failure mode of Fisher-based
arguments~{\cite{vallisneri2008fisher}}, so we type every use explicitly
and adhere to the typing throughout.

\paragraph{U1 --- Local information geometry (exact).}
As the pullback of the waveform-space metric, $\Gamma(f)$ and its
eigenstructure are exact local properties of the likelihood at the
expansion point, independent of the posterior's global shape. The
misalignment $C(f)$, the commutator \eqnref{commutator_def}, and their
evolution with cutoff are U1 statements and require no Gaussianity
assumption. Because local geometry can vary along an extended ridge,
all U1 quantities are computed at three expansion points spanning the
empirically probed ridge segment, and the spread is reported as a
systematic band [\manifest{T6: ridge-uniformity band}].

\paragraph{U2 --- Global covariance proxy (approximate; its failure is
quantified, not assumed away).}
The Laplace approximation
$\rho_f \approx \mathcal{N}(x_{*}(f), \Gamma(f)^{-1})$ is invalid along
$\vd$, where the true posterior is a prior-bounded ridge. We do not
rely on it there; instead we \emph{measure} its failure with the
non-Gaussianity index
\begin{equation}
N(f) \;=\;
\frac{\Var_{\rho_f}\!\bigl(\vd^{\top} x\bigr)}
{\bigl(\Gamma(f)^{-1}\bigr)_{\vd\vd}} ,
\label{eq:ngauss}
\end{equation}
the ratio of the true posterior variance along the degeneracy direction
to the Fisher prediction; $N \approx 1$ certifies the Gaussian regime
and $N \gg 1$ quantifies the breakdown. The single-pass filter bias of
\secref{exactness} accrues precisely where $N$ is large, so U2 failure
is not a caveat to this work but part of its subject matter.

\paragraph{U3 --- Pairwise discrimination (quadratic order; exact
upgrade available).}
The scalar $\rhosq(f)$ of \eqnref{rhosq} approximates the discriminating
signal-to-noise between two specified points to second order in the
displacement. Where the displacement is not small --- as for the ridge
endpoints of GW170817 --- we compute the exact mismatch directly from
two waveform evaluations and quote the quadratic form only alongside
its exact value [\manifest{T5: quadratic-to-exact ratio}].

Finally, we verify the typing empirically with the consistency
criterion of Ref.~{\cite{vallisneri2008fisher}}: at each cutoff and along
each eigendirection $v_i$, displace by one Fisher standard deviation
$\lambda_i^{-1/2}$ and compare the exact mismatch to its quadratic
prediction. The resulting pass/fail map over $(f, v_i)$
(\figref{validity}) shows the linear-signal regime holding along $v_1$
throughout the band and failing along $\vd$ at low cutoffs, with the
onset of failure tracking $N(f)$ --- a direct empirical delineation of
where each role of $\Gamma$ is licensed. Numerical conditioning of the
eigenproblem (the smallest spectral gaps, and machine-precision
derivatives in place of finite differences) is documented in
Appendix~{\ref{sec:appendix}}.